\section{Introduction}
\label{sec:intro}

Most communication technology runs on spoken and written language and serves sign language
users poorly \citep{bragg2019interdisciplinary,yin2021including}. A recognizer that runs in a
browser on an ordinary laptop, without sending video anywhere, could support practice tools and
simple aids for a fixed vocabulary. Landmark-based recognizers suit this setting: MediaPipe
keypoints \citep{lugaresi2019mediapipe} are small, largely appearance-free and fast to classify. The
Google ISLR data \citep{kaggle2023islr} provides them for \V{data.signs} isolated signs, collected
with the partners behind PopSign, a game that helps hearing parents of deaf children learn to sign
\citep{starner2023popsign}; the vocabulary is everyday and child-oriented.

Test accuracy, however, says little about the same model on a webcam. The live pipeline has to
reproduce, in another language and with another tracker, every step that produced the training
features: which hand goes in which slot, how coordinates
are normalized, how time is resampled, where a sign starts and ends. Any mismatch is train--serve
skew \citep{sculley2015debt,breck2017mltestscore}, and offline evaluation cannot see it. Our
first browser demo shipped a model with \V{m.cnn.top1}\% test top-1 and recognized
\V{lg.web.c0headbuild.top1}\% of new webcam clips.

We measure that gap for ManuVox, our in-browser recognizer, on a low-power laptop. Our
contributions are:
\begin{itemize}
\item A signer-independent comparison of landmark models with paired, Holm-corrected tests
  (Section~\ref{sec:models}): a convolution--transformer hybrid and a longer recipe improve test
  top-1 by \V{cmp.ship.diff} points over a CNN, for each of the \V{data.testsigners} test signers
  (weak evidence about new signers), while richer input features alone did not help measurably. On an Intel i3-1115G4 laptop the
  hybrid classifies a sign in \V{b.wasm1.p50} ms in Chrome (Section~\ref{sec:efficiency}).
\item An audit of six candidate train--serve discrepancies, with an ablation of their fixes on
  \V{lg.clips} new webcam clips through two pipelines (Section~\ref{sec:livegap}). Assigning hands
  to slots as the training tracker did has by far the largest effect (\V{lg.py.c2hands.dtop1}
  points with the training tracker, \V{lg.web.c2hands.dtop1} in the browser); the other changes
  are smaller or inconclusive with fifty clips from one signer. The fixes recover
  \V{lg.web.recover}--\V{lg.py.recover}\% of the hybrid's drop from its test accuracy, so most of
  the gap remains.
\item A cross-dataset test and a second training set. On the webcam recordings of ASL Citizen's
  unseen test signers, v2 reaches \V{cit.v2.top1}\% on \V{cit.signs} mapped signs, and its
  calibration degrades. v3, trained also on ASL Citizen and with webcam augmentation, reaches
  \V{cit.v3.top1}\% there, an in-domain score, and gains \V{v3c.ks.diff} points on the Google ISLR
  test split. Live, it gains \V{v3c.web.diff} points of clip top-1 in the browser and
  \V{v3c.py.diff} in Python, neither significant with fifty clips (Section~\ref{sec:liveresults}).
\item A personalization feature, Teach Dex, evaluated separately (Section~\ref{sec:teach}):
  examples replayed from one recording per sign raise clip top-1 on a second take from
  \V{td.v2.model}\% to \V{td.v2.teach}\% with v2, in the best case of one signer and one sitting.
  A v3 run of the same test was compromised by its harness.
\end{itemize}
One feature specification shared by Python and the browser, and a golden-fixture test between
them, now guard against the two pipelines drifting apart (Section~\ref{sec:faults}).
